\mxfront{stmg-front}
\mxcover{MATHÉMATIQUES}{Terminale STMG}{Livret annuel de cours et d'entraînement}{55 séances structurées • 14 sujets de DS}{2026--2027}
\chapter*{Mode d'emploi}
\addcontentsline{toc}{chapter}{Mode d'emploi}
Ce livret suit les 55 séances prévues du 15 septembre 2026 au 20 mai 2027. Chaque chapitre débute par un cours de référence illustré. Chaque fiche de séance contient un rituel d'automatismes, un rappel ciblé, une activité de recherche, la trace écrite attendue, un exemple rédigé, des exercices gradués, au moins un exercice d'application original, un billet de sortie et un travail personnel.

\begin{itemize}
\item En \textbf{séance d'une heure}, conserver les automatismes, la trace écrite, l'exemple et les exercices marqués essentiels.
\item En \textbf{séance de deux heures}, mener l'activité de construction, ajouter les exercices de consolidation et terminer par le billet de sortie.
\item Les exercices d'application sont originaux. Les encadrés « Pour prolonger » renvoient aux pages et aux exercices du manuel \emph{Calao Terminale STMG} pour s'entraîner davantage ; aucun énoncé du manuel n'est reproduit.
\item Les corrections, les anticipations d'erreurs et des indications de conduite de classe apparaissent uniquement dans l'édition enseignant.
\end{itemize}

\begin{keypoint}
Le livret articule les attendus du programme de mathématiques de terminale technologique publié au BO spécial n° 8 du 25 juillet 2019 et la progression annuelle fournie. Les notions sont réactivées régulièrement par les automatismes et les séances de consolidation.
\end{keypoint}

\tableofcontents
